\section{Современные методы регистрации нейронной активности in vivo}
\label{sec:registration_methods}

Изучение коллективной динамики нейронных популяций невозможно без экспериментальных методов, позволяющих одновременно регистрировать активность большого числа нервных клеток в мозге бодрствующих и свободно ведущих себя животных. За последние десятилетия число одновременно записываемых нейронов увеличивалось экспоненциально~\cite{Stevenson2011}, что принципиально изменило возможности анализа популяционной активности и поиска скрытых переменных в нейронных данных. В данном разделе рассматриваются основные классы методов регистрации нейронной активности in~vivo: электрофизиологические и оптические.

\subsection{Электрофизиологические методы}
\label{sec:electrophysiology}

Электрофизиологические методы регистрации основаны на измерении электрических потенциалов, возникающих при генерации потенциалов действия (спайков) нервными клетками. Исторически именно внеклеточная электрофизиология положила начало изучению нейронной активности на уровне отдельных клеток и их популяций.

Пионерские работы Хьюбела и Визеля по исследованию рецептивных полей нейронов зрительной коры~\cite{HubelWiesel1962} были выполнены с помощью единичных микроэлектродов. Этот подход позволял регистрировать активность одного или нескольких нейронов с высоким временным разрешением (порядка миллисекунд), однако принципиально ограничивал размер регистрируемой популяции. Открытие клеток места в гиппокампе О'Кифом и Достровским~\cite{o1971hippocampus} также было получено с помощью одиночных электродов для внеклеточной записи у свободно двигающихся крыс.

Переход к мультиэлектродным массивам позволил существенно увеличить число одновременно записываемых нейронов. Тетроды~--- пучки из четырёх электродов~--- дали возможность более надёжно разделять активность нескольких нейронов на основе форм волн спайков~\cite{Wilson1993}. Планарные кремниевые зонды и массивы электродов (типа Utah array) расширили масштаб записей до сотен каналов, однако подходы на основе таких массивов были ограничены числом доступных каналов усиления и объёмом записываемых данных.

Качественный скачок в масштабах электрофизиологических записей произошёл с появлением нейропиксельных зондов (Neuropixels), содержащих до нескольких тысяч электродов на одном кремниевом зонде~\cite{Hong2019}. Эти зонды позволяют одновременно регистрировать активность сотен и даже тысяч нейронов из множества структур мозга с миллисекундным временным разрешением. Высокая плотность электродов обеспечивает возможность качественной сортировки спайков и идентификации отдельных нервных клеток. Технологии нового поколения продолжают развиваться в направлении увеличения числа каналов и снижения инвазивности~\cite{Hong2019}.

Ключевыми достоинствами электрофизиологических методов остаются высокое временное разрешение и прямая регистрация электрической активности нейронов. Вместе с тем, они имеют ряд ограничений: инвазивность, трудности с идентификацией типов клеток и отслеживанием одних и тех же нейронов на протяжении длительных экспериментов (дней и недель).

\subsection{Оптические методы}
\label{sec:optical_methods}

Развитие оптических методов регистрации нейронной активности открыло принципиально новые возможности для исследования коллективной динамики нейронных популяций. В основе этих методов лежит визуализация внутриклеточной динамики ионов кальция (Ca$^{2+}$), концентрация которых возрастает при генерации потенциалов действия.

\paragraph{Генетически кодируемые кальциевые индикаторы.}
Визуализация кальциевой активности стала возможной благодаря разработке генетически кодируемых кальциевых индикаторов (GECI, genetically encoded calcium indicators). GECI представляют собой белки, флуоресценция которых изменяется при связывании ионов Ca$^{2+}$. Наиболее широко используемым семейством GECI является GCaMP~--- химерный белок, состоящий из зелёного флуоресцентного белка (GFP), кальмодулина и пептида M13, конформация которого меняется при связывании Ca$^{2+}$, что приводит к изменению интенсивности флуоресценции. Разработка улучшенных вариантов GCaMP (GCaMP5, GCaMP6)~\cite{Chen2013} обеспечила высокую чувствительность, позволяющую детектировать одиночные потенциалы действия. Новейшие индикаторы серии jGCaMP8~\cite{Zhang2023} обладают ещё более быстрой кинетикой (постоянная времени спада --- десятки миллисекунд), что приближает временное разрешение оптических записей к электрофизиологическому.

Важным преимуществом генетической экспрессии индикаторов является возможность нацеливания на определённые типы клеток с помощью специфических промоторов, а также стабильная экспрессия на протяжении недель и месяцев, что позволяет отслеживать одни и те же нейроны в хронических экспериментах.

\paragraph{Двухфотонная микроскопия.}
Двухфотонная лазерная сканирующая микроскопия является одним из основных инструментов визуализации кальциевой активности нейронных популяций. Метод основан на нелинейном возбуждении флуоресценции двумя фотонами инфракрасного диапазона, что обеспечивает оптическое сечение~--- возбуждение флуоресценции строго в фокальной плоскости~--- и глубокое проникновение в ткань мозга (до нескольких сотен микрометров). Двухфотонная микроскопия обеспечивает субклеточное пространственное разрешение и позволяет одновременно визуализировать активность сотен и тысяч нейронов~\cite{Pachitariu2017}. Недавние разработки, в том числе двухфотонные минископы, позволили проводить крупномасштабную двухфотонную визуализацию кальциевой активности у свободно двигающихся мышей~\cite{Zong2022}.

\paragraph{Минископия.}
Принципиальным шагом к изучению нейронных популяций у свободно ведущих себя животных стала минископия~--- визуализация кальциевой активности с помощью миниатюрных головных флуоресцентных микроскопов (минископов). В сочетании с имплантируемым оптическим зондом (градиентной линзой, GRIN) минископ позволяет записывать активность нейронных популяций в глубоких структурах мозга, таких как гиппокамп, у свободно двигающихся животных. Важным достоинством минископии является возможность хронической регистрации одних и тех же нейронов на протяжении дней и недель~\cite{Sheintuch2017}, что делает этот метод незаменимым инструментом для изучения динамики нейронных ансамблей в ходе обучения и формирования памяти~\cite{Sotskov2022}.

\paragraph{Особенности кальциевого имиджинга.}
При всех достоинствах оптических методов следует учитывать их принципиальные отличия от электрофизиологии. Кальциевый сигнал является непрямым индикатором электрической активности: флуоресценция отражает динамику внутриклеточного Ca$^{2+}$, а не спайки непосредственно. Временное разрешение кальциевого имиджинга существенно уступает электрофизиологическому~--- даже быстрейшие индикаторы имеют постоянную времени спада порядка десятков миллисекунд~\cite{Zhang2023}, что приводит к временной свёртке сигнала. Было показано, что динамика популяции нейронов, измеренная с помощью кальциевого имиджинга, может количественно отличаться от электрофизиологических данных~\cite{Wei2020}. Тем не менее на популяционном уровне основные характеристики коллективной активности~--- низкоразмерные траектории, корреляционная структура~--- воспроизводятся обоими методами~\cite{Wei2020}.

В целом, развитие оптических методов регистрации, особенно в сочетании с генетически кодируемыми индикаторами и миниатюрными микроскопами, создало экспериментальную основу для крупномасштабного исследования коллективной динамики нейронных популяций, которое является предметом настоящей диссертации.
